%%%%%%%%%%%%%%%%%%%%%%%%%%%%%%%%%%%%%%%%%
% Long Lined Cover Letter
% LaTeX Template
% Version 1.0 (1/6/13)
%
% This template has been downloaded from:
% http://www.LaTeXTemplates.com
%
% Original author:
% Matthew J. Miller
% http://www.matthewjmiller.net/howtos/customized-cover-letter-scripts/
%
% License:
% CC BY-NC-SA 3.0 (http://creativecommons.org/licenses/by-nc-sa/3.0/)
%
%%%%%%%%%%%%%%%%%%%%%%%%%%%%%%%%%%%%%%%%%

%----------------------------------------------------------------------------------------
%	PACKAGES AND OTHER DOCUMENT CONFIGURATIONS
%----------------------------------------------------------------------------------------

\documentclass[10pt,stdletter,dateno,sigleft,letterpaper]{newlfm} % Extra options: 'sigleft' for a left-aligned signature, 'stdletternofrom' to remove the from address, 'letterpaper' for US letter paper - consult the newlfm class manual for more options

\usepackage{charter} % Use the Charter font for the document text
\usepackage{epstopdf} % Use eps images even for PDF outputs

%Put 1.5 line spacing for corrections before making the letter final.
\usepackage{setspace}
\singlespacing
%\onehalfspacing
%\doublespacing

\usepackage[pdfsubject={Cover Letter for Autodesk},pdfauthor={Jorge Garcia},pdftitle={Cover Letter - Jorge Garcia},pdfkeywords={Jorge A. Garcia, Cover letter, Autodesk}]{hyperref}

\newsavebox{\Luiuc}\sbox{\Luiuc}{\parbox[b]{1.75in}{\vspace{0.25in}
\includegraphics[width=1.0\linewidth]{img/PolyLogo.eps}}} % Company/institution logo at the top left of the page
\makeletterhead{Uiuc}{\Lheader{\usebox{\Luiuc}}}

\newlfmP{sigsize=8pt} % Slightly decrease the height of the signature field
\newlfmP{addrfromphone} % Print a phone number under the sender's address
\newlfmP{addrfromemail} % Print an email address under the sender's address
\PhrPhone{Phone} % Customize the "Telephone" text
\PhrEmail{Email} % Customize the "E-mail" text

\lthUiuc % Print the company/institution logo




%----------------------------------------------------------------------------------------
%	YOUR NAME AND CONTACT INFORMATION
%----------------------------------------------------------------------------------------

\namefrom{Jorge A. Garc\'{i}a Galicia} % Name

\addrfrom{
\today\\[12pt] % Date
510 Salem Street, Apt 102 \\ % Address
Lafayette, IN 47904
}

\phonefrom{(765) 543-9834} % Phone number

\emailfrom{garci191@purdue.edu} % Email address

%----------------------------------------------------------------------------------------
%	ADDRESSEE AND GREETING/CLOSING
%----------------------------------------------------------------------------------------

\greetto{Dear Hiring manager,} % Greeting text

\closeline{
Sincerely, \newline
	\parbox[b]{0.6in}{
		\includegraphics[width=1.2\linewidth]{img/jorgeFirma.eps}
	} 
} % Closing text


%\nameto{Mrs. Jane Smith} % Addressee of the letter above the to address

\addrto{Autodesk, Inc.\\
  111 McInnis Parkway,\\
  San Rafael, CA 94903
}

%----------------------------------------------------------------------------------------

\begin{document}
\begin{newlfm}

%----------------------------------------------------------------------------------------
%	LETTER CONTENT
%----------------------------------------------------------------------------------------


%The first paragraph should contain a self-introduction. You should write who you are, where your expertise lies, where you found the job posting, and why you want to apply for the job.
%I am a PhD candidate with a strong background in mathematics and currently focusing on computer science. I found out about the Lead Principal Engineer position for Choreograph Cooperating 3D Printing Robots at Autodesk from LinkedIn. I believe I am an excellent candidate for this job: first, my research has been focused on additive manufacturing and second, my professional expertise is in software development. Since I expect to graduate in August 2017, I am looking to start a career in technology and Autodesk is the perfect place to do so.
I am a Ph.D. candidate with a strong background in mathematics and currently focusing on Computer Science. I found out about the Principal Software Engineer position for Rendering in A360 at Autodesk from LinkedIn. I believe I am an excellent candidate for this job: first, my graduate studies have been in Computer Graphics, since my Master's and second, my professional expertise is also in software development on graphics, moreover my last internship (Summer 2016) was about creating a rendering engine. Since I expect to graduate in August 2017, I am looking to start a career in software development focused in Computer Graphics and Autodesk is the perfect place to do so.


%The second paragraph should respond directly to the job description written by the hiring manager. Describe how your previous job experiences, skills, and abilities will allow you to meet the company’s needs. To make that easier, you can (and should) literally include words and phrases from the job description in your cover letter.
%As a student in Computer Science I am a proficient programmer, especially in C++, C\# and Java, and because my focus has been in computer graphics I have experience using OpenGL and Qt. Even though most of my experience is in desktop applications, I have excellent web developing skills. During the last academic year I was a teaching assistant in two web development classes at the Computer Graphics department at Purdue. One of this classes was focused in front end technologies using HTML5, CSS3 and JavaScript and the second was focused in back end development using ASP.NET MVC. I am also familiar with PHP programming and frameworks like Joomla and Wordpress.
%As a student in Computer Science, I am a proficient programmer and algorithmic thinker, especially in C/C++ and because my focus has been in computer graphics, I have experience using OpenGL and Qt. Even though I have been focusing on software developing during the last 10 years, I have an excellent background in Computational Geometry and Geometry Processing. While earning my master's in Computer Science I took two classes focused on Computational Geometry, and I have always worked on research in Geometry Processing. During fall of 2017, I took a digital manufacturing class, where I used CNC machines and a laser cutter along with 3D printers.
Beside Computer Sciences, I also have background in Applied Mathematics, because of that I am a proficient programmer and algorithmic thinker especially in C/C++ and because of computer graphics I also have experience using OpenGL and Qt for creating user applications. During my PhD, part of my research involved parallelization on GPU using CUDA and compute shaders. And during my last internship, I worked in a deferred rendering engine for a research group, where I implemented most of the algorithms in GLSL. Since then, I have been working in improvement in my OpenGL skills. In previous years I also did two internships at the most important GPU manufacturer company in the world; there, I learn about control version and automatic testing, so I am more than ready to work in a production environment.

%To go the extra mile, do some research about the company, and try to find out what they are doing — and why — given the current state of their industry. Explain how you can fit into that schema, and help push the company forward and achieve any goals you suspect they may have.

During my doctoral studies, I become familiar with several Autodesk products. I wrote a research article in 3D printing were I compare my algorithm for creating support structures to the one used in Meshmixer. I also have experience using Maya to model and render 3D meshes, for both my research work and my teaching. Finally, even when my professional background is technical, I come from a family of artists, I have experience doing live coding and I even coauthored an art paper at SIGGRAPH so I am very positive about working in a company were artist and scientist work in close collaboration.

%During my doctoral studies, I coauthored a research article in 3D printing where I compare my algorithm for creating support structures to the one used in Meshmixer. I also coauthored an article about toolpath generation in SIGGRAPH. During the last three summers I have held internships where I was developing production code, so I am very familiar with the release cycle for software, including testing. Finally, even though my professional background is technical, I come from a family of artists. I have experience doing live coding and I even coauthored an art paper at SIGGRAPH.

%The final paragraph is called the “call to action” portion of your cover letter. Inform them that you’d love to get interviewed. Give them your contact information. Tell them that you’ll be in contact with them in a week if you don’t hear back. Thank them for spending the time to read your cover letter.
I would love to have an interview regarding the Principal Software Engineer position for Rendering in A360. My preferred contact method is email at garci191@purdue.edu. Thank you very much the time you took reading this letter and your consideration. I am eagerly looking forward for your response.
%I would love to have an interview regarding the Lead Principal Engineer position for Choreograph Cooperating 3D Printing Robots. My preferred contact method is email at garci191@purdue.edu. Thank you very much for the time you took reading this letter and your consideration. I am eagerly looking forward for your response.

%----------------------------------------------------------------------------------------

\end{newlfm}
\end{document}